% Ad M. Brutum 2.3 --- headnote
% ad-brutum-02-03, 1 April 43 BC, Dyrrachium

\textit{M. Junius Brutus to Cicero, written from Dyrrachium
on the kalends of April 43 BC --- Perseus dateline
\textit{Scr.\ Dyrrach.\ i K.\ Apr.\ a.\ 711 (43)}, confirmed
by the closing subscript \textit{Kalend.\ Apr.\ Dyrrhachio}.
The meta entry carries the impossible date
\textit{-0043-03-32} (32 March); the Perseus dateline
fixes it as 1 April 43 BC, and the corrected date is used
in the parallel sidecar. Brutus is writing from
Dyrrachium, where he has consolidated his hold on the
Balkan provinces after taking over Hortensius's legions
and capturing C. Antonius, Mark Antony's brother, at
Apollonia. News of the murder of Trebonius by Dolabella
at Smyrna has reached him; news of Mutina has not.}

\textit{The voice is Brutus's: short clauses, plain
verbs, the soldier's directness. The six sections divide
into three pairs of practical questions. Section 1
acknowledges the loss of Asia along with Trebonius, with
the lapidary observation that the province can be
recovered but a recovery after the fact is no less
disgraceful for being possible. Section 2 raises the
matter that will dominate this correspondence: what to do
with the captive C.\ Antonius, whom Brutus is plainly
considering sparing --- ``his entreaties move me'' ---
against the wishes of others (the \textit{aliquorum furor}
of section 2 is the senatorial party in Rome and Cicero
himself, who wants him executed). Section 3 reports
Cassius's seizure of Syria with the Syrian legions
under Murcus and Marcius, and notes the discretion
Brutus has used in containing the news at the family
level until he hears Cicero's view. Section 4 is the
warm and slightly teasing acknowledgement of Cicero's
two recent speeches (Philippic V on the kalends of
January, and a defence against Calenus): Brutus concedes
the joke and grants Cicero the title \textit{Philippicae}
for his consular orations. Section 5 turns to logistics
--- money and supplement --- and to the grief of Asia's
loss. Section 6 closes with the unaffected praise of
Cicero's son, then on Brutus's staff: a father's son,
who will earn his honours on his own footing. The
letter is one of the most characteristic short
specimens of Brutus's prose --- terse, candid,
unornamented, written by a commander to a statesman
whom he treats as an equal.}
